\documentclass[10pt,a4paper]{amsart}
\usepackage[margin=19mm]{geometry}
\usepackage{amsmath,amssymb,mathtools}
\usepackage[scr=rsfs,cal=euler]{mathalfa}
\usepackage{xfrac}
\usepackage[unicode,hidelinks,bookmarks=false]{hyperref}
\newcommand{\ie}{i.e.\ }
\newcommand{\cf}{cf.\ }
\newcommand{\resp}{respectively\ }
\newcommand{\Subsection}{\subsection}
\providecommand{\nobreakdash}{\nobreak\hbox{-}}
\setlength{\emergencystretch}{2em}
\multlinegap0pt
\usepackage{xcolor}
\newtheorem{proposition}{Proposition}
\newtheorem{lemma}{Lemma}
\newtheorem{theorem}{Theorem}
\def\calC{\mathcal{C}}
\def\mod{\text{mod}\;}
\title{long version here}
\author{Olivia Dumitrescu, Rick Miranda}
\date{Source: arXiv:2205.13605v1 (CC BY 4.0)}
\begin{document}
\maketitle

\noindent\emph{Source and licence.} long version here. Version arXiv:2205.13605v1 (CC BY 4.0), distributed by arXiv under the Creative Commons Attribution 4.0 International licence, http://creativecommons.org/licenses/by/4.0/, as recorded in the arXiv licence metadata for this exact version and shown on the abstract page https://arxiv.org/abs/2205.13605v1. Full licence notice: \texttt{licenses/arxiv-2205.13605-CC-BY-4.0.txt}.\par\smallskip

\noindent\textit{Editorial scope.} Counted unit: the article's theorem thm:0weyllines (source0.tex lines 1821--1831). The article's complete original proof of it is delivered with it (source0.tex lines 1833--1917, one \texttt{proof} environment of 85 lines). No mathematics is written, repaired or re-derived: the statement and the proof are the article's own lines, in the article's own notation, and no retained line is substituted or re-flowed.  Retained uncounted as the source-local context the proof needs: lines 964--970; lines 1265--1267; lines 1338--1344.  Nothing was omitted from any retained span, so the package carries no omission record.  No retained line contains a citation, so the wrapper carries no bibliography and no bibliography entry.  No retained line contains a figure environment, an image-inclusion command or drawing code, so no image is delivered and the package declares no asset dependency.  Editorial formatting only, in addition: an AMS \texttt{amsart} preamble that restates the article's own declarations and macros used by the retained lines, namely the environments \texttt{proposition}, \texttt{lemma}, \texttt{theorem} on separate counters, together with the standard packages those lines need.  The retained environments print in this document as Proposition 1, Lemma 1, Theorem 1 in order of appearance, whereas the article prints the same items with the numbering of its own full document.  The complete native source of the article as distributed in the arXiv e-print for this exact version is bundled unchanged as \texttt{sources/126206/source0.tex}.
\begin{proposition}\label{CXprop}
Suppose $X$ and $X'$ are subsets of the basis vectors $\{x_i\}$,
and $w$ is in the Weyl group $W$.
If $w(\mathcal{C}_X) \cap \mathcal{C}_{X'} \neq \emptyset$,
then $X=X'$,  $w \in W_X$, and $w(\mathcal{C}_X) = \mathcal{C}_{X}$.
In particular the stabilizer of any $f \in \mathcal{C}_X$ is $W_X$.
\end{proposition}

\begin{equation}\label{finiteWeylrs}
(r+1)^2 > s(r-1) \iff r=2, s \leq 8; r=3, s \leq 7; r=4, s \leq 8; r\geq 5, s \leq r+3.
\end{equation}

\begin{lemma}\label{twoCremonareducedclasses}
Suppose that $c$ and $c'$ are two Cremona reduced classes
with multiplicities in non-increasing order.
Assume that they are Cremona-equivalent:
there is an element $w$ in the Weyl group such that $w(c)=c'$.
Then $c \equiv c' \mod F$.
\end{lemma}

\begin{theorem}\label{thm:0weyllines}
Fix $r\geq 3$.
Then the only Cremona-reduced classes $c = (d;\underline{m})$
(with all $d$, $m_i$ non-negative)
which are Cremona-equivalent to $h-e_1$
are the classes $h-e_i$ for $1 \leq i \leq s$.
Any $(0)$-Weyl line with positive degree and non-negative multiplicities
may be brought to $h-e_1$ via a series of Cremona transformations,
each of which monotonically decreases the degree,
followed by a single permutation.
\end{theorem}

\begin{proof}
Take $c$ and permute the multiplicities to be in non-increasing order.
Then both $h-e_1 \mod F$ and $c\mod F$
are in the closure $\bar\calC$ of the positive chamber, and so we may apply
Lemma \ref{twoCremonareducedclasses}
and conclude (since they are Cremona equivalent)
that they are equal $\mod F$.
Hence we may write $c = aF + (h-e_1)$ for some integers $a \geq 0$.

Since $\langle c,F\rangle = \langle h-e_1,F\rangle = 2$,
this forces $a\langle F,F\rangle = 0$.
If $\langle F, F \rangle \neq 0$, we conclude that $a=0$.

Suppose that $\langle F, F \rangle = 0$.
Then $2-r = \langle c,c\rangle = \langle h-e_1,h-e_1\rangle + 2a\langle h-e_1,F\rangle
= 2-r +4a$
and we conclude $a=0$ in this case as well.

Hence $c = h-e_1$ after permuting the multiplicities,
and so $c = h-e_j$ for some $j$ as claimed.

The final statement follows from applying the Cremona reduction algorithm.

%The statement is equivalent to asserting that the only Cremona-reduced class
%with all $d$ and $m_i$ non-negative and with the multiplicities in descending order
%is the class $h-e_1$.
%
%Suppose $c$ is another such class.  
%Then both $h-e_1(\mod F)$ and $c(\mod F)$
%are in the closure $\bar{\mathcal{C}}$ of the positive chamber.
%We see that $h-e_1 (\mod F)$ is in $\mathcal{C}_X$ 
%for the subset $X = \{X_0, X_i\;|\;2\leq i\leq s-1\}$.
%Therefore since $c$ and $h-e_1$ are Cremona-equivalent, there
%is an element of the Weyl group that takes $c$ to $h-e_1$;
%by Proposition \ref{CXprop}, 
%we have that $c(\mod F)$ is in $\mathcal{C}_X$ also.
%This implies that if we write $c = dh-\sum_i m_i$ with $m_i$ in descending order,
%we have $d = \sum_{i=1}^{r+1}m_i$, $m_1>m_2$ and all other $m_i$ equal to $m_2$.
%We may therefore write $c(\mod F) = dh-m_1e_1-\sum_{i=2}^s m_2 e_i$
%for non-negative integers $d$, $m_1$, and $m_2$
%satisfying those equalities and inequalities.
%We conclude that this is also true for the class $c$; adding multiples of $F$ do not alter those inequalities or equalities.
%Hence we may write $c = aF + b(h-e_1)$ for integers $a \geq 0$, $b \geq 1$
%(here $a = m_2$, $b = m_1-m_2$, and $d = m_1+rm_2 = (r+1)a+b$).
%
%The Weyl group preserves the intersection form 
%so we also have $\langle c,F\rangle = \langle h-e_1,F\rangle = 2-r$,
%which gives that $a\langle F,F\rangle + 2b = 2-r$.
%
%Since $r \geq 3$, this quantity is negative.
%In the case that $\langle F, F \rangle \geq 0$,
%this forces $b=0$, which is a contradiction.
%Hence we may focus on the non-Mori-Dream-Space case where $\langle F, F \rangle < 0$.
%
%If $r$ is odd, then $\langle F, F \rangle$ is even,
%and therefore $a\langle F,F\rangle + 2b = 2-r$ is impossible.
%Hence we may assume $r \geq 4$ and is even.
%
%Consider now $\langle c, c \rangle$; since $c$ is Cremona equivalent to $h=e_1$, this quantity is $2$; expanding it out using $c = aF + b(h-e_1)$ then gives the equation
%\[
%2 = a^2\langle F, F\rangle + 4ab + (2-r)b^2
%\]
%This is a quadratic form in the variables $a,b$, and since $\langle F, F \rangle < 0$,
%we will have a contradiction when this form is negative semi-definite.
%This will happen when the determinant is non-negative, i.e., when
%$\langle F, F \rangle (2-r) \geq 4$.
%This is automatic if $r \geq 6$, so we only have to consider the case $r=4$ now.
%
%When $r=4$, this condition fails only when $\langle F, F \rangle =-1$,
%and this is not possible for any $s$.
%
%Hence
%\begin{align*}
%2 &= (r+1)d-(r-1)m_1-(r-1)(s-1)m_2 \\
%&= (r+1)(m_1+rm_2) -(r-1)m_1-(r-1)(s-1)m_2 \\
%&= 2m_1 + (r^2+2r-1-(r-1)s)m_2 \\
%&\geq (r^2+2r-1-(r-1)s +2)m_2 + 2 \;\;\text{since}\;\; m_1 \geq m_2+1)\\
%&= ((r+1)^2-(r-1)s)m_2+2 \\
%&> 2 \;\;\text{if}\;\; m_2 > 0,
%\end{align*}
%the last strict inequality coming from (\ref{finiteWeylrs}).
%This gives a contradiction unless $m_2=0$.
%Hence $c = m_1(h-e_1)$, and $c\cdot K_Y = 2$ forces $m_1=1$ as required.

\end{proof}

\end{document}
